\documentclass[a4paper,14pt]{extarticle}

\usepackage[a4paper,margin=2.5cm]{geometry}
\usepackage{iftex}
\ifPDFTeX
  \usepackage[T2A]{fontenc}
  \usepackage[utf8]{inputenc}
  \usepackage[russian]{babel}
  \usepackage{tempora}
\else
  \usepackage[russian]{babel}
  \usepackage{fontspec}
  \setmainfont{Times New Roman}
\fi

\usepackage{amsmath,amssymb}
\usepackage{booktabs}
\usepackage{longtable}
\usepackage{array}
\usepackage{setspace}
\usepackage[unicode,colorlinks=true,linkcolor=blue,urlcolor=magenta]{hyperref}

\onehalfspacing
\renewcommand{\arraystretch}{1.15}

\newcommand{\FormulaNote}[2]{%
  \begin{equation*}
  #1
  \end{equation*}
  \noindent\textbf{Пример применения.} #2\par\medskip
}

\title{Математическая статистика\\Формулы и примеры применения}
\author{Конспект по материалам \texttt{matstat.pdf}}
\date{}

\begin{document}
\maketitle
\tableofcontents
\newpage

\section{Вариационные ряды и эмпирическое распределение}

\subsection{Частоты и интервалы}

\FormulaNote{w_i=\frac{m_i}{n}}{Переходит от абсолютной частоты к относительной; показывает долю наблюдений с вариантом $x_i$.}

\FormulaNote{\sum_{i=1}^k m_i=n \qquad \sum_{i=1}^k w_i=1}{Проверяет полноту таблицы частот и нормировку относительных частот.}

\FormulaNote{k\approx 1+3{,}322\,\lg n}{Используется для выбора числа интервалов при построении интервального вариационного ряда.}

\FormulaNote{h=\frac{x_{\max}-x_{\min}}{k}}{Задаёт ширину интервала после того, как число интервалов уже выбрано.}

\subsection{Эмпирическая функция распределения}

\FormulaNote{F_n^*(x)=\frac{1}{n}\sum_{i=1}^{n} I\{X_i\le x\}}{Строит эмпирическую функцию распределения, то есть долю наблюдений, не превосходящих $x$.}

\FormulaNote{F_n^*(x)=\frac{m_x}{n}}{Удобная запись ЭФР через число элементов выборки, которые лежат левее или на границе $x$.}

\section{Выборочные характеристики}

\subsection{Среднее арифметическое}

\FormulaNote{\bar{x}=\frac{1}{n}\sum_{i=1}^{n}x_i}{Основная оценка центра выборки; показывает средний уровень признака.}

\FormulaNote{\bar{x}=\frac{\sum_{i=1}^{k}x_i^*m_i}{\sum_{i=1}^{k}m_i}=\frac{\sum_{i=1}^{k}x_i^*m_i}{n}}{Та же формула для вариационного ряда с частотами; используется, когда данные уже сгруппированы.}

\FormulaNote{\sum_{i=1}^{n}(x_i-\bar{x})=0}{Показывает, что положительные и отрицательные отклонения от среднего взаимно компенсируются.}

\FormulaNote{\sum_{i=1}^{k}m_i(x_i^*-\bar{x})=0}{То же свойство для данных, записанных с частотами.}

\subsection{Мода и медиана}

\FormulaNote{\text{Mo}=x_j,\qquad m_j=\max_i m_i}{Для дискретного ряда находит наиболее часто встречающееся значение признака.}

\FormulaNote{\text{Mo}=x_0+\frac{m_0-m_{-1}}{(m_0-m_{-1})+(m_0-m_{+1})}\,h}{Приближённо вычисляет моду для интервального ряда по модальному интервалу.}

\FormulaNote{\text{Me}=
\begin{cases}
x_{(n+1)/2}, & n \text{ нечётно},\\[0.5ex]
\dfrac{x_{(n/2)}+x_{(n/2+1)}}{2}, & n \text{ чётно}.
\end{cases}}{Находит медиану дискретного упорядоченного ряда.}

\FormulaNote{\text{Me}=x_0+\frac{0{,}5n-n_{\text{Me}-1}}{m_{\text{Me}}}\,h}{Приближённо вычисляет медиану для интервального ряда.}

\subsection{Разброс и вариация}

\FormulaNote{R=x_{\max}-x_{\min}}{Даёт самый простой показатель разброса: расстояние между крайними значениями выборки.}

\FormulaNote{d=\frac{1}{n}\sum_{i=1}^{n}|x_i-\bar{x}|}{Считает среднее линейное отклонение от среднего.}

\FormulaNote{d=\frac{\sum_{i=1}^{k}m_i|x_i^*-\bar{x}|}{\sum_{i=1}^{k}m_i}}{Та же характеристика для данных, сгруппированных по значениям.}

\FormulaNote{D=\frac{1}{n}\sum_{i=1}^{n}(x_i-\bar{x})^2}{Выборочная дисперсия: измеряет средний квадрат отклонений от среднего.}

\FormulaNote{D=\frac{\sum_{i=1}^{k}m_i(x_i^*-\bar{x})^2}{\sum_{i=1}^{k}m_i}}{Формула дисперсии для вариационного ряда с частотами.}

\FormulaNote{D=\overline{x^2}-(\bar{x})^2,\qquad \overline{x^2}=\frac{1}{n}\sum_{i=1}^{n}x_i^2}{Удобная вычислительная формула для дисперсии; обычно быстрее считать через квадраты.}

\FormulaNote{\sigma_v=\sqrt{D}}{Среднее квадратическое отклонение, то есть корень из дисперсии; измеряется в тех же единицах, что и исходные данные.}

\FormulaNote{S^2=\frac{n}{n-1}D=\frac{1}{n-1}\sum_{i=1}^{n}(x_i-\bar{x})^2}{Исправленная выборочная дисперсия; применяется как несмещённая оценка генеральной дисперсии.}

\FormulaNote{S=\sqrt{S^2}}{Исправленное среднее квадратическое отклонение; используется, когда нужна несмещённая оценка разброса.}

\FormulaNote{d_Q=\frac{Q_3-Q_1}{2}}{Квартильное отклонение: показывает разброс центральной половины данных.}

\FormulaNote{V=\frac{S}{\bar{x}}\cdot 100\%}{Коэффициент вариации; позволяет сравнивать относительную изменчивость разных выборок.}

\section{Выборочные моменты и свойства оценок}

\subsection{Моменты}

\FormulaNote{\nu_r=\frac{1}{n}\sum_{i=1}^{n}x_i^r}{Начальный выборочный момент порядка $r$; используется для описания формы распределения и в методе моментов.}

\FormulaNote{\nu_r=\frac{\sum_{i=1}^{k}m_i(x_i^*)^r}{\sum_{i=1}^{k}m_i}}{Та же формула для сгруппированных данных с частотами.}

\FormulaNote{\mu_r=\frac{1}{n}\sum_{i=1}^{n}(x_i-\bar{x})^r}{Центральный выборочный момент порядка $r$; нужен для оценки разброса, асимметрии и эксцесса.}

\FormulaNote{\mu_r=\frac{\sum_{i=1}^{k}m_i(x_i^*-\bar{x})^r}{\sum_{i=1}^{k}m_i}}{Центральный момент для вариационного ряда с частотами.}

\FormulaNote{\mu_2=D}{Связывает второй центральный момент с дисперсией.}

\subsection{Свойства оценок}

\FormulaNote{\mathbb{E}(\hat{\theta})=\theta}{Условие несмещённости оценки.}

\FormulaNote{\hat{\theta}\xrightarrow{P}\theta}{Условие состоятельности оценки: при росте объёма выборки оценка сходится к истинному параметру.}

\FormulaNote{\mathbb{E}\bar{X}=a}{Показывает, что выборочное среднее является несмещённой оценкой математического ожидания.}

\FormulaNote{D\bar{X}=\frac{\sigma^2}{n}}{Показывает, что разброс выборочного среднего уменьшается при росте объёма выборки.}

\FormulaNote{\mathbb{E}D=\frac{n-1}{n}\sigma^2}{Показывает смещение обычной выборочной дисперсии.}

\FormulaNote{\mathbb{E}S^2=\sigma^2}{Доказывает, что исправленная дисперсия является несмещённой оценкой генеральной дисперсии.}

\subsection{Групповая и общая дисперсии}

\FormulaNote{\bar{x}=\frac{\sum_{j}N_j\bar{x}_j}{N}}{Находит общее среднее по группам, если совокупность предварительно разбита на классы.}

\FormulaNote{D_{\text{гр}}=\frac{\sum_{j}N_jD_j}{N}}{Средняя из внутригрупповых дисперсий; показывает разброс внутри групп.}

\FormulaNote{D_{\text{меж}}=\frac{\sum_{j}N_j(\bar{x}_j-\bar{x})^2}{N}}{Межгрупповая дисперсия; отражает различия между средними групп.}

\FormulaNote{D_{\text{общ}}=\frac{\sum_{j}\sum_{i}(x_{ij}-\bar{x})^2}{N}}{Общая дисперсия для всей совокупности.}

\FormulaNote{D_{\text{общ}}=D_{\text{гр}}+D_{\text{меж}}}{Правило сложения дисперсий: общая изменчивость раскладывается на внутригрупповую и межгрупповую части.}

\section{Методы точечного оценивания}

\subsection{Метод моментов}

\FormulaNote{
\begin{cases}
\mathbb{E}_{\theta}X=\nu_1,\\
\mathbb{E}_{\theta}X^2=\nu_2,\\
\vdots\\
\mathbb{E}_{\theta}X^m=\nu_m
\end{cases}
}{Система метода моментов: теоретические моменты приравниваются к выборочным, чтобы получить оценки параметров.}

\FormulaNote{\mathbb{E}X=\frac{\theta}{2},\qquad \hat{\theta}=2\bar{X}}{Пример метода моментов для равномерного распределения $U(0,\theta)$ по первому моменту.}

\FormulaNote{\mathbb{E}X^2=\frac{\theta^2}{3},\qquad \hat{\theta}=\sqrt{3\,\overline{X^2}}}{Пример метода моментов для $U(0,\theta)$ по второму моменту.}

\FormulaNote{\mathbb{E}X=\frac{1}{\lambda},\qquad \hat{\lambda}=\frac{1}{\bar{X}}}{Пример метода моментов для экспоненциального распределения.}

\subsection{Метод максимального правдоподобия}

\FormulaNote{L(\theta)=\prod_{i=1}^{n}f(x_i,\theta)}{Функция правдоподобия: показывает, насколько хорошо параметр $\theta$ объясняет наблюдённую выборку.}

\FormulaNote{\hat{\theta}=\arg\max_{\theta\in\Theta}L(\theta)}{Определение оценки максимального правдоподобия.}

\FormulaNote{\ell(\theta)=\ln L(\theta)=\sum_{i=1}^{n}\ln f(x_i,\theta)}{Логарифм правдоподобия; его обычно максимизируют вместо самой функции правдоподобия.}

\FormulaNote{\frac{d\ell(\theta)}{d\theta}=0}{Уравнение, которое решают для поиска точки максимума, если параметр один.}

\FormulaNote{\frac{\partial \ell}{\partial \theta_j}=0,\qquad j=1,\dots,m}{Система условий максимума для случая нескольких параметров.}

\FormulaNote{f(x,\theta)=\frac{1}{\theta},\qquad 0\le x\le \theta}{Плотность равномерного распределения $U(0,\theta)$, нужная для вывода МП-оценки.}

\FormulaNote{L(\theta)=\frac{1}{\theta^n}I\{\theta\ge x_{(n)}\}}{Правдоподобие для $U(0,\theta)$; используется для получения оценки $\hat{\theta}=x_{(n)}$.}

\FormulaNote{\hat{\theta}=x_{(n)}}{Оценка максимального правдоподобия для параметра $\theta$ в равномерном распределении $U(0,\theta)$.}

\section{Доверительные интервалы}

\subsection{Общие определения}

\FormulaNote{P(\theta_1<\theta<\theta_2)=\gamma}{Определяет смысл доверительного интервала: параметр попадает в интервал с вероятностью $\gamma$.}

\FormulaNote{\alpha=1-\gamma}{Связывает доверительную вероятность и уровень значимости.}

\subsection{Среднее при известной дисперсии}

\FormulaNote{Z=\frac{\bar{X}-a}{\sigma/\sqrt{n}}\sim N(0,1)}{Переход к стандартному нормальному закону при известной дисперсии.}

\FormulaNote{\bar{X}-u_\gamma\frac{\sigma}{\sqrt{n}}<a<\bar{X}+u_\gamma\frac{\sigma}{\sqrt{n}}}{Доверительный интервал для математического ожидания при известной дисперсии.}

\FormulaNote{n\ge \frac{u_\gamma^2\sigma^2}{\Delta^2}}{Формула для оценки требуемого объёма выборки при заданной предельной ошибке.}

\subsection{Среднее при неизвестной дисперсии}

\FormulaNote{t=\frac{\bar{X}-a}{S/\sqrt{n}}}{Статистика Стьюдента, когда дисперсия неизвестна и заменяется исправленным стандартным отклонением.}

\FormulaNote{\bar{X}-t_{\alpha/2,n-1}\frac{S}{\sqrt{n}}<a<\bar{X}+t_{\alpha/2,n-1}\frac{S}{\sqrt{n}}}{Доверительный интервал для математического ожидания при неизвестной дисперсии.}

\subsection{Доля}

\FormulaNote{\hat{p}=\frac{m}{n},\qquad \hat{q}=1-\hat{p}}{Выборочная доля и дополнение к ней; используются при построении интервала для вероятности события.}

\FormulaNote{\hat{p}-u_\gamma\sqrt{\frac{\hat{p}\hat{q}}{n}}<p<\hat{p}+u_\gamma\sqrt{\frac{\hat{p}\hat{q}}{n}}}{Приближённый доверительный интервал для доли в большой выборке.}

\FormulaNote{\Delta=u_\gamma\sqrt{\frac{\hat{p}\hat{q}}{n}}}{Формула предельной ошибки для оценки доли.}

\FormulaNote{n\ge \frac{u_\gamma^2pq}{\Delta^2}}{Оценка необходимого объёма выборки, если заранее известны $p$ и $q$.}

\FormulaNote{n\ge \frac{u_\gamma^2}{4\Delta^2}}{Худший случай для оценки объёма выборки, когда $p$ заранее неизвестно.}

\subsection{Дисперсия}

\FormulaNote{\chi^2=\frac{(n-1)S^2}{\sigma^2}}{Статистика хи-квадрат для нормальной совокупности; используется при построении интервала для дисперсии.}

\FormulaNote{\frac{(n-1)S^2}{\chi^2_{1-\alpha/2,n-1}}<\sigma^2<\frac{(n-1)S^2}{\chi^2_{\alpha/2,n-1}}}{Доверительный интервал для генеральной дисперсии.}

\section{Статистические гипотезы и критерии проверки}

\subsection{Нулевая и альтернативная гипотезы}

\FormulaNote{H_0:\ a=a_0 \qquad H_1:\ a\neq a_0}{Типичный вид гипотезы о среднем в двусторонней постановке.}

\FormulaNote{H_0:\ \sigma_X^2=\sigma_Y^2 \qquad H_1:\ \sigma_X^2\neq \sigma_Y^2}{Гипотеза о равенстве дисперсий двух совокупностей.}

\subsection{Проверка среднего}

\FormulaNote{Z=\frac{\bar{X}-a_0}{\sigma/\sqrt{n}}}{Статистика для проверки гипотезы о среднем, когда дисперсия известна.}

\FormulaNote{|Z_{\text{набл}}|>u_{1-\alpha/2}}{Правило отклонения нулевой гипотезы для двустороннего z-критерия.}

\FormulaNote{t=\frac{\bar{X}-a_0}{S/\sqrt{n}}}{Статистика для проверки среднего при неизвестной дисперсии.}

\FormulaNote{|t_{\text{набл}}|>t_{\alpha/2,n-1}}{Правило отклонения нулевой гипотезы для двустороннего t-критерия.}

\FormulaNote{
t=\frac{\bar{X}-\bar{Y}}{S_p\sqrt{\frac{1}{n_1}+\frac{1}{n_2}}},
\qquad
S_p^2=\frac{(n_1-1)S_X^2+(n_2-1)S_Y^2}{n_1+n_2-2}
}{Критерий для сравнения двух независимых средних при равных, но неизвестных дисперсиях.}

\FormulaNote{k=n_1+n_2-2}{Число степеней свободы в двухвыборочном t-критерии.}

\FormulaNote{d_i=x_i-y_i,\qquad t=\frac{\bar{d}}{S_d/\sqrt{n}}}{Парный t-критерий: сначала считают разности, затем проверяют гипотезу о нулевом среднем разностей.}

\subsection{Проверка дисперсий и согласия}

\FormulaNote{F=\frac{S_X^2}{S_Y^2}}{Статистика критерия Фишера--Снедекора для проверки равенства дисперсий.}

\FormulaNote{\chi^2_{\text{набл}}=\sum_{i=1}^{k}\frac{(n_i-np_i)^2}{np_i}}{Критерий согласия Пирсона: сравнивает наблюдаемые и теоретические частоты.}

\FormulaNote{\nu=k-1}{Число степеней свободы в критерии Пирсона, если параметры распределения не оценивались по выборке.}

\FormulaNote{\nu=k-m-1}{Число степеней свободы в критерии Пирсона, если по выборке оценивалось $m$ параметров.}

\FormulaNote{\hat{n}_{ij}=\frac{n_{i\cdot}n_{\cdot j}}{n}}{Ожидаемая частота в таблице сопряжённости при условии независимости признаков.}

\FormulaNote{\chi^2_{\text{набл}}=\sum_{i=1}^{r}\sum_{j=1}^{c}\frac{(n_{ij}-\hat{n}_{ij})^2}{\hat{n}_{ij}}}{Критерий независимости признаков по таблице сопряжённости.}

\FormulaNote{\nu=(r-1)(c-1)}{Число степеней свободы в критерии независимости.}

\section{Корреляционный анализ}

\subsection{Коэффициент Пирсона}

\FormulaNote{
r_{xy}=
\frac{\sum_{i=1}^{n}(x_i-\bar{x})(y_i-\bar{y})}
{\sqrt{\sum_{i=1}^{n}(x_i-\bar{x})^2}\,\sqrt{\sum_{i=1}^{n}(y_i-\bar{y})^2}}
}{Измеряет тесноту линейной связи между двумя количественными признаками.}

\FormulaNote{r_{xy}=\frac{\overline{xy}-\bar{x}\bar{y}}{\sigma_x\sigma_y}}{Эквивалентная вычислительная форма коэффициента корреляции Пирсона.}

\FormulaNote{t_{\text{набл}}=\frac{r\sqrt{n-2}}{\sqrt{1-r^2}}}{Статистика для проверки значимости коэффициента корреляции.}

\FormulaNote{|t_{\text{набл}}|>t_{\alpha/2,n-2}}{Правило отклонения гипотезы об отсутствии линейной корреляции.}

\subsection{Ранговая и согласованность}

\FormulaNote{r_s=1-\frac{6\sum_{i=1}^{n}d_i^2}{n(n^2-1)}}{Коэффициент ранговой корреляции Спирмена; применяется для ранговых данных.}

\FormulaNote{t=\frac{r_s\sqrt{n-2}}{\sqrt{1-r_s^2}}}{Приближённая проверка значимости ранговой корреляции Спирмена.}

\FormulaNote{S=\sum_{i=1}^{n}(R_i-\bar{R})^2}{Промежуточная сумма квадратов отклонений рангов в коэффициенте конкордации Кендалла.}

\FormulaNote{W=\frac{12S}{m^2(n^3-n)}}{Коэффициент конкордации Кендалла; измеряет согласованность мнений экспертов.}

\FormulaNote{\chi^2_{\text{набл}}=m(n-1)W}{Статистика для проверки значимости коэффициента конкордации.}

\section{Регрессионный анализ}

\subsection{Парная регрессия}

\FormulaNote{y=f(x)+\varepsilon}{Общая запись регрессии: наблюдаемое значение состоит из систематической части и случайной ошибки.}

\FormulaNote{y=a+bx+\varepsilon}{Линейная парная регрессия; подходит, когда связь между признаками близка к прямой.}

\FormulaNote{Q(a,b)=\sum_{i=1}^{n}(y_i-a-bx_i)^2\to\min}{Метод наименьших квадратов: параметры выбирают так, чтобы сумма квадратов ошибок была минимальной.}

\FormulaNote{
\begin{cases}
na+b\sum x_i=\sum y_i,\\
a\sum x_i+b\sum x_i^2=\sum x_i y_i
\end{cases}
}{Нормальные уравнения метода наименьших квадратов для линейной регрессии.}

\FormulaNote{b=\frac{\overline{xy}-\bar{x}\bar{y}}{\overline{x^2}-(\bar{x})^2},\qquad a=\bar{y}-b\bar{x}}{Явная формула для оценок коэффициентов линейной регрессии.}

\FormulaNote{r_{xy}=\frac{b\sigma_x}{\sigma_y}}{Связывает наклон линейной регрессии с коэффициентом корреляции.}

\FormulaNote{R^2=\frac{\sum(\hat{y}_i-\bar{y})^2}{\sum(y_i-\bar{y})^2}}{Коэффициент детерминации: показывает долю объяснённой вариации результата.}

\FormulaNote{R^2=1-\frac{\sum(y_i-\hat{y}_i)^2}{\sum(y_i-\bar{y})^2}}{Эквивалентная форма коэффициента детерминации через остатки.}

\FormulaNote{R^2=r_{xy}^2}{Для парной линейной регрессии коэффициент детерминации равен квадрату корреляции.}

\FormulaNote{F_{\text{набл}}=\frac{R^2}{1-R^2}(n-2)}{Критерий Фишера для проверки значимости парной линейной регрессии.}

\FormulaNote{F_{\text{набл}}=\frac{R^2/p}{(1-R^2)/(n-p-1)}}{Критерий Фишера для множественной линейной регрессии.}

\FormulaNote{F_{\text{набл}}>F_{\alpha,p,n-p-1}}{Правило отклонения нулевой гипотезы о незначимости регрессии.}

\subsection{Линеаризация нелинейных моделей}

\FormulaNote{
\begin{array}{lcll}
y=a+bx+cx^2 &\Rightarrow& z=x^2, & y=a+bx+cz,\\[0.5ex]
y=ab^x      &\Rightarrow& \ln y=\ln a+x\ln b, & Y=A+Bx,\\[0.5ex]
y=ax^b      &\Rightarrow& \ln y=\ln a+b\ln x, & Y=A+BX,\\[0.5ex]
y=ae^{bx}   &\Rightarrow& \ln y=\ln a+bx, & Y=A+bx,\\[0.5ex]
y=a+b\ln x  &\Rightarrow& X=\ln x, & y=a+bX
\end{array}
}{Эта схема показывает, как нелинейную зависимость сводят к линейной заменой переменных и затем решают обычную задачу МНК.}

\subsection{Множественная регрессия}

\FormulaNote{y=b_0+b_1x_1+b_2x_2+\dots+b_px_p+\varepsilon}{Множественная линейная регрессия: результат зависит сразу от нескольких факторов.}

\FormulaNote{\mathbf{y}=\mathbf{X}\mathbf{b}+\varepsilon}{Матричная форма записи множественной регрессии.}

\FormulaNote{\hat{\mathbf{b}}=(\mathbf{X}^T\mathbf{X})^{-1}\mathbf{X}^T\mathbf{y}}{Матричная формула МНК для оценки коэффициентов множественной регрессии.}

\FormulaNote{
\begin{cases}
nb_0+b_1\sum x_1+b_2\sum x_2=\sum y,\\
b_0\sum x_1+b_1\sum x_1^2+b_2\sum x_1x_2=\sum x_1y,\\
b_0\sum x_2+b_1\sum x_1x_2+b_2\sum x_2^2=\sum x_2y
\end{cases}
}{Нормальные уравнения для регрессии с двумя факторами.}

\section*{Краткий итог}

Формулы из \texttt{matstat.pdf} используются для трёх основных задач:
\begin{itemize}
  \item описать выборку и её разброс;
  \item оценить неизвестные параметры генеральной совокупности;
  \item проверить гипотезы и построить модели зависимости между признаками.
\end{itemize}

\end{document}
